\documentclass[conference,letterpaper]{IEEEtran}
\IEEEoverridecommandlockouts

\usepackage[T1]{fontenc}
\usepackage{amsmath,amssymb}
\usepackage{booktabs}
\usepackage{cite}
\usepackage{newtxtext,newtxmath}
\usepackage{xcolor}
\usepackage{tikz}
\usepackage{url}
\usetikzlibrary{arrows.meta,positioning}

\makeatletter
\newcommand{\linebreakand}{%
  \end{@IEEEauthorhalign}
  \hfill\mbox{}\par
  \mbox{}\hfill\begin{@IEEEauthorhalign}}
\makeatother

\newcommand{\modelStats}{Gaze-summary MLP}
\newcommand{\modelTime}{Dwell-time LR}
\newcommand{\modelLinear}{Gaze-summary LR}
\newcommand{\modelGru}{Gaze GRU}
\newcommand{\modelTeacher}{Text--gaze teacher}
\newcommand{\modelText}{Text-only control}
\newcommand{\modelDistill}{Distilled GRU}
\newcommand{\best}[1]{\textbf{#1}}

\title{Compact Gaze-Only Reading-Goal Classification\\
with Word-Aligned Teacher Distillation}

\author{
\IEEEauthorblockN{Hyunsu Go}
\IEEEauthorblockA{Seoul National University\\
Seoul, Republic of Korea\\
hsmail02@snu.ac.kr}
\and
\IEEEauthorblockN{Sumin Lee}
\IEEEauthorblockA{Seoul National University\\
Seoul, Republic of Korea\\
cirtuare@snu.ac.kr}
\and
\IEEEauthorblockN{Jinyong Jun}
\IEEEauthorblockA{Seoul National University\\
Seoul, Republic of Korea\\
junjinyong@snu.ac.kr}
\linebreakand
\IEEEauthorblockN{Kyeonghun Kim}
\IEEEauthorblockA{OUTTA\\
Republic of Korea\\
kyeonghun.kim@outta.ai}
\and
\IEEEauthorblockN{Hyuk-Jae Lee}
\IEEEauthorblockA{Seoul National University\\
Seoul, Republic of Korea\\
hyukjae@snu.ac.kr}
\and
\IEEEauthorblockN{Nam-Joon Kim\textsuperscript{*}%
\thanks{\textsuperscript{*}Corresponding author: Nam-Joon Kim.}}
\IEEEauthorblockA{Seoul National University\\
Seoul, Republic of Korea\\
knj01@snu.ac.kr}}

\begin{document}
\maketitle

\begin{abstract}
Reading interfaces could benefit from inferring whether a user is reading for
general comprehension or seeking specific information.  Prior multimodal
systems retain passage text, a language encoder, and fixation-to-word
alignment at inference.  We instead train a compact gaze-only GRU with a
word-aligned teacher used only during training.  On OneStop, under simultaneous
reader and article holdout, distillation raises balanced accuracy by 3.39
percentage points over the same GRU trained with hard labels.  Reader-clustered
and crossed reader/article bootstrap intervals exclude zero.  Label smoothing
does not reproduce the gain, and within-class shuffling of teacher logits
reduces it by 1.90 points.  Removing contextual text from the teacher has no
measurable effect on student performance.  The distilled model requires only
fixation sequences at inference and retains the baseline architecture and
cost.  Across observed sequence lengths, its model-only latency is 5.08~ms at
the median and 16.57~ms at the 95th percentile on a single-thread laptop CPU.
\end{abstract}

\begin{IEEEkeywords}
eye tracking, scanpath analysis, reading behavior, human--computer interaction,
privileged information
\end{IEEEkeywords}

\section{Introduction}

A reader may approach the same passage to understand it broadly or to locate
specific information.  These goals are associated with differences in
fixation durations, regressions, and scanpath structure
~\cite{rayner1998,shubi2023information}.  An opt-in reading-goal classifier
could support navigation or content assistance after a paragraph
~\cite{shubi2025decoding}.  Such a classifier must generalize to
readers and articles absent from its training data.

Recent work established paragraph-level reading-goal decoding on OneStop Eye
Movements and found that models combining scanpaths with contextual language
representations achieve the strongest performance~\cite{shubi2025decoding}.
Such models require passage text, fixation-to-word alignment, and a language
encoder at inference.  We ask whether these resources can instead be confined
to training while prediction uses gaze alone.

We study a text--gaze teacher that combines contextual word representations
with word-level gaze summaries and a compact student that observes the
chronological fixation sequence.  The primary comparison holds the student
architecture and input fixed.  At inference, the student uses no teacher, text
encoder, tokenization, or fixation-to-word alignment.

Because each OneStop reader completes only one goal condition and articles
recur across readers, random trial splits can place the same reader or article
in both training and test data.  We therefore hold out readers and articles
simultaneously.  Uncertainty is estimated by resampling readers and, in a
sensitivity analysis, readers and articles together.

This study makes three contributions:

\begin{itemize}
  \item We compare distilled and undistilled gaze-only GRUs with identical
  inputs, architecture, training data, and evaluation trials.
  \item We compare logistic gaze baselines and probe the gain with label
  smoothing, within-class target shuffling, and contextual-text removal.
  \item We quantify the parameter count, serialized size, and single-thread
  CPU latency across the observed sequence-length distribution.
\end{itemize}

\begin{figure}[t]
\centering
\definecolor{trainblue}{RGB}{49,92,145}
\definecolor{infergreen}{RGB}{49,120,88}
\begin{tikzpicture}[
  x=1mm,
  y=1mm,
  font=\scriptsize,
  module/.style={draw, rounded corners=2pt, align=center, inner sep=3pt,
                 minimum width=28mm, minimum height=7mm},
  trainmodule/.style={module, fill=trainblue!8, draw=trainblue},
  infermodule/.style={module, fill=infergreen!10, draw=infergreen},
  trainarr/.style={-{Latex[length=1.6mm]}, thick, draw=trainblue},
  inferarr/.style={-{Latex[length=1.6mm]}, thick, draw=infergreen},
  scanarr/.style={-{Latex[length=1.1mm]}, line width=.35pt,
                  draw=infergreen!80!black},
  kdarr/.style={-{Latex[length=1.6mm]}, thick, dashed,
                draw=gray!75!black}]
  \node[font=\scriptsize\bfseries, text=black, anchor=west]
        at (0,24) {Training only};
  \begin{scope}[shift={(3,15)}]
    \draw[gray!70, line width=.45pt] (0,3) -- (22,3);
    \draw[gray!70, line width=.45pt] (0,0) -- (18,0);
    \draw[gray!70, line width=.45pt] (0,-3) -- (21,-3);
    \fill[trainblue] (4,3) circle (.75);
    \fill[trainblue] (11,3) circle (1.05);
    \fill[trainblue] (16,0) circle (.65);
    \fill[trainblue] (7,0) circle (.90);
    \fill[trainblue] (3,-3) circle (.70);
    \fill[trainblue] (14,-3) circle (1.00);
    \node[text=black, anchor=north] at (11,-4.5)
          {Text + word-level gaze};
  \end{scope}
  \node[trainmodule, text=black] (teacher) at (46,15)
        {Text--gaze teacher};
  \draw[trainarr] (25,15) -- (teacher.west);

  \begin{scope}[shift={(70,15)}]
    \fill[trainblue!45] (0,.8) rectangle (4,1.6);
    \fill[trainblue] (0,-1.6) rectangle (10,-.8);
    \node[text=black, anchor=south] at (5,3.2)
          {Teacher prediction};
  \end{scope}
  \draw[trainarr] (teacher.east) -- (70,15);

  \node[font=\scriptsize\bfseries, text=black, anchor=west]
        at (0,3) {Inference path};
  \begin{scope}[shift={(3,-5)}]
    \coordinate (s1) at (1,2);
    \coordinate (s2) at (10,2);
    \coordinate (s3) at (20,2);
    \coordinate (turnR) at (20,0);
    \coordinate (turnL) at (2,0);
    \coordinate (s4) at (2,-2);
    \coordinate (s5) at (17,-2);
    \draw[scanarr] (s1) -- (s3);
    \draw[scanarr] (s3) -- (turnR) -- (turnL) -- (s4);
    \draw[scanarr] (s4) -- (s5);
    \fill[infergreen] (s1) circle (.65);
    \fill[infergreen] (s2) circle (.90);
    \fill[infergreen] (s3) circle (.55);
    \fill[infergreen] (s4) circle (.90);
    \fill[infergreen] (s5) circle (.65);
    \node[text=black, anchor=north] at (11,-3.5)
          {Fixation sequence};
  \end{scope}
  \node[infermodule, text=black] (student) at (46,-5)
        {Gaze-only student};
  \draw[inferarr] (25,-5) -- (student.west);

  \begin{scope}[shift={(70,-5)}]
    \fill[infergreen!45] (0,.8) rectangle (4.8,1.6);
    \fill[infergreen] (0,-1.6) rectangle (9.2,-.8);
    \node[text=black, anchor=north] at (5,-2.7)
          {Student prediction};
  \end{scope}
  \draw[inferarr] (student.east) -- (70,-5);

  \draw[kdarr] (75,12) -- (75,-2);
\end{tikzpicture}
\caption{Training-only distillation.  The teacher uses contextual text and
word-level gaze summaries.  The student uses a fixation sequence.  The
simplified scanpath retains line progression and a return sweep from a
representative trial, with marker size indicating fixation duration.  The
dashed link denotes soft-target matching.  Only the green path remains at
inference.}
\label{fig:overview}
\end{figure}

\section{Related Work}

\subsection{Eye Movements and Reading Goals}

OneStop includes eye movements from ordinary reading for comprehension,
information seeking, and repeated reading~\cite{berzak2025onestop}.  It has
also supported gaze-only and text--gaze models for predicting reading
comprehension under new-reader and new-item conditions
~\cite{shubi2024comprehension}.

Shubi \emph{et al.} introduced the reading-goal decoding task used here and
evaluated gaze-only and multimodal models under new-reader and new-item
conditions~\cite{shubi2025decoding}.  That study compares model families and
partial scanpaths, and its highest-performing model retains text at inference.  We
instead hold the gaze-only inference architecture fixed and test whether it
benefits from richer training-time supervision.

\subsection{Distillation with Privileged Information}

Knowledge distillation (KD) trains a student with class labels and a teacher's
temperature-softened predictive distribution~\cite{hinton2015distilling}.
Generalized distillation allows the teacher to use privileged information that
is unavailable to the student at inference~\cite{lopezpaz2016unifying}.  Here,
the teacher receives contextual word representations and word-level gaze
summaries, whereas the student receives fixation features.  The gaze-only
teacher control removes the contextual word representations while preserving
the teacher architecture and gaze aggregation.

\section{Method}

\subsection{Task and Data}

The unit of analysis is a completed paragraph-reading trial, represented by its
eye movements.  The binary target distinguishes ordinary reading for
comprehension ($y=0$) from information seeking ($y=1$).  We analyze the OneStop
Eye Movements corpus~\cite{berzak2025onestop} and retain first readings from
these two conditions.  Repeated readings, question-screen events, answers,
correctness, and all post-reading information are excluded.

OneStop was collected with participant consent under an institutional ethics
protocol and released publicly after anonymization~\cite{berzak2025onestop}.
Our analysis uses only this public release.

The resulting data contain 19,438 trials from 360 readers and 30 articles,
with nearly balanced goal labels.  Each reader completed one goal condition,
and the median trial contains 96 valid fixations.  Reader identity, article
identity, question content, session position, and condition metadata are not
model inputs.

Each fixation is encoded by ten features: $x$ and $y$ coordinates normalized
to the paragraph bounds; log fixation duration; horizontal and vertical
displacement from the preceding fixation; log saccade amplitude; sine and
cosine of the saccade angle; a line-transition indicator; and sequence position
within the observed sequence.  Retained fixations have positive finite durations
and finite paragraph-screen coordinates.  Continuous features are standardized
with statistics estimated only from the current training set.

\subsection{Models}

Table~\ref{tab:models} summarizes the five neural model families.  Models
trained without distillation minimize binary cross-entropy (BCE).  \modelStats{} maps
22 sequence summaries to a 128--64 multilayer perceptron with ReLU and dropout
0.2.  The summaries cover duration, displacement, saccade amplitude, line
transitions, spatial dispersion, and path length.

\begin{table}[t]
\caption{Model objectives, inputs, and trainable parameter counts.}
\label{tab:models}
\centering
\scriptsize
\setlength{\tabcolsep}{3.4pt}
\begin{tabular}{@{}lllr@{}}
\toprule
Model & Objective & Model input & Parameters \\
\midrule
\modelStats{} & BCE & gaze summaries & 11,265 \\
\modelGru{} & BCE & fixation sequence & 43,797 \\
\modelTeacher{} & BCE & text + word-level gaze & 1,038,945$^a$ \\
\modelText{} & BCE & text & 1,038,945$^a$ \\
\modelDistill{} & BCE + KD & fixation sequence & 43,797 \\
\bottomrule
\multicolumn{4}{@{}l}{\scriptsize $^a$Trainable parameters; frozen RoBERTa is excluded.}
\end{tabular}
\end{table}

\modelGru{} applies layer normalization to each fixation and processes the
sequence with a two-layer unidirectional GRU~\cite{cho2014gru} with hidden size
64.  A 64--64--1 ReLU head with dropout 0.2 produces the logit.  This model is the
undistilled comparison for \modelDistill{}.  Sequences are zero-padded within
each minibatch.  The classifier uses the GRU output at the trial's final
fixation before padding.

A frozen RoBERTa-base encoder produces subword
representations~\cite{liu2019roberta}.  \modelTeacher{} averages the
representations belonging to each paragraph word.  Eight word-level gaze
features encode fixation count, total and mean dwell time, first-fixation and
first-pass dwell time, regression-in and revisit counts, and skipping.  Text
and gaze features are projected to 128 and 64 dimensions, concatenated with a
32-dimensional learned word-position embedding, and mapped to width 192.  A
learned classification token and a three-layer, six-head Transformer
~\cite{vaswani2017attention} with feed-forward width 384, GELU, and dropout 0.1
produce the logit.  A padding mask excludes padded word positions.
\modelText{} retains the text branch and fusion architecture but sets all gaze
features to zero.  It tests whether paragraph text alone predicts the reading
goal.

\modelDistill{} has exactly the same architecture and inputs as \modelGru{}.
Let $z_S$ and $z_T$ be the student and teacher logits, and let $\tau=2$.  The
student minimizes
\begin{equation}
\begin{aligned}
\mathcal{L}_D={}&(1-\alpha)\,\mathrm{BCE}(y,\sigma(z_S))\\
&+\alpha\tau^2\,\mathrm{KL}\!\left[
q_T^{(\tau)}\,\middle\|\,q_S^{(\tau)}\right],
\quad q_j^{(\tau)}=\mathrm{Bern}\!\left(\sigma\!\left(\frac{z_j}{\tau}\right)\right).
\end{aligned}
\label{eq:kd}
\end{equation}
Teacher logits are generated only for student-training trials within each outer
split.  No held-out test trial is used to fit either model or construct a
distillation target.

\subsection{Evaluation Protocol}

Readers are divided into five goal-stratified folds, and articles are divided
into five folds.  Their Cartesian product gives 25 disjoint test partitions
that together cover all retained trials once.  For each reader--article fold
pair, training excludes every trial from either held-out fold, and the test set
is their intersection.  All test readers and articles are therefore absent
from training.

Within each outer training set, model selection uses three paired inner
validation splits.  Each validation split likewise holds out both its readers
and articles from the corresponding inner training set.  Learning rates are
selected from
$\{10^{-3},3\!\times\!10^{-4}\}$ for the compact models and
$\{3\!\times\!10^{-4},10^{-4}\}$ for the teacher models.  We also select
$\alpha\in\{0.25,0.50,0.75\}$ for \modelDistill{}.  Balanced accuracy is the
selection metric, with negative log-likelihood (NLL) as the first tie-breaker.
Each final model is refitted for the median number of epochs selected across
the three inner splits.  The decision threshold is fixed at 0.5 before test
evaluation.  Test labels are not used for early stopping, calibration, or
threshold selection.

Models are optimized with AdamW~\cite{loshchilov2019adamw}, weight decay
$10^{-4}$, and gradient clipping at 1.0.  Inner training uses at most 30 epochs
with a patience of five epochs.  Batch sizes are 512 for \modelStats{}, 256 for
the GRU students, and 128 for the teacher models.  Final refits use the selected
learning rate and fixed epoch count described above, without test-dependent
stopping.

\subsection{Post Hoc Ablations}

After the primary result, we ran post hoc controls using the same outer splits
and training setup.  Paired conditions inherit the selected learning rate and
epoch count and keep the remaining student, teacher, or KD settings fixed.

The first control replaces teacher targets with binary label smoothing
($\epsilon=0.10$)~\cite{szegedy2016label}.  The second zeros every RoBERTa
representation in an architecture-matched teacher while retaining its
word-level gaze summaries.  Across three fixed random seeds, we refit the BCE
and label-smoothed GRUs, both teachers, and both KD students.  Each student
uses logits from the teacher trained with the same seed.

The third control cyclically reassigns teacher logits within each ground-truth
class, preserving the class-conditional distributions while removing
trial-level correspondence.  It uses the same seeds and student settings.

We also fit logistic controls using either log dwell time per
word~\cite{shubi2025decoding} or the 22 gaze summaries.  Their regularization
is selected within the same nested reader/article splits.

\subsection{Metrics and Uncertainty}

Balanced accuracy (BA) is the primary metric because it averages the two
goal-specific recalls~\cite{brodersen2010balanced}.  Secondary metrics are
AUROC, macro F1, negative log-likelihood (NLL), Brier score, and 15-bin expected
calibration error (ECE)~\cite{guo2017calibration}.

For the primary contrast between \modelDistill{} and \modelGru{}, a paired
5,000-draw percentile bootstrap resamples whole readers within each goal.  We
report the central 95\% interval.  A 2,000-draw sensitivity analysis resamples
readers within each goal and independently resamples articles, weighting each
trial by the product of their multiplicities.
Bootstrap samples are paired between models.  We also report the number
of outer partitions in which the primary contrast is positive.  These
partitions are not treated as independent trials.  For each post hoc contrast,
resampled units are paired across conditions, and the seed-specific BA
differences are averaged within each draw.

\subsection{CPU Measurement}

For timing, we select the distilled development checkpoint whose full-sequence
validation BA is closest to the mean across all inner evaluations.  After 100
warm-up iterations, we time one full-sequence forward pass for each retained
trial in randomized order.  Measurements use an Intel Core Ultra 7 155H and
one PyTorch thread, excluding input construction, parsing, and disk I/O.

\begin{table}[t]
\caption{Pooled out-of-fold performance.  Lower is better for NLL, Brier, and
ECE.  LR controls were added post hoc.}
\label{tab:main-results}
\centering
\scriptsize
\setlength{\tabcolsep}{2.1pt}
\begin{tabular}{@{}lcccccc@{}}
\toprule
Model & BA & AUROC & Macro F1 & NLL & Brier & ECE \\
\midrule
\modelTime{}$^b$ & .5887 & .6292 & .5887 & .6635 & .2359 & .0155 \\
\modelLinear{}$^b$ & .6066 & .6493 & .6065 & .6572 & .2327 & .0236 \\
\modelStats{} & .5942 & .6359 & .5936 & .6632 & .2360 & .0410 \\
\modelGru{} & .5866 & .6178 & .5862 & .6710 & .2391 & .0210 \\
\modelTeacher{} & .6126 & .6613 & .6120 & .6561 & .2312 & .0496 \\
\modelText{} & .4986 & .5017 & .4986 & .6952 & .2510 & .0304 \\
\modelDistill{} & \best{.6205} & \best{.6708} & \best{.6204} & \best{.6427} &
\best{.2265} & \best{.0142} \\
\bottomrule
\multicolumn{7}{@{}l}{\scriptsize $^b$Nested selection within the frozen reader/article splits.}
\end{tabular}
\end{table}

\section{Results}

\subsection{Confirmatory Performance}

Table~\ref{tab:main-results} reports pooled out-of-fold predictions.  \modelDistill{}
improves BA over \modelGru{} by 0.0339.  The reader-clustered 95\% interval is
$[0.0175, 0.0504]$, and the crossed reader/article interval is
$[0.0130, 0.0553]$.  Both exclude zero.

The improvement is positive in 18 of the 25 test partitions.  \modelDistill{} also exceeds
\modelStats{} and \modelTeacher{} in pooled BA, whereas \modelText{} remains at
chance.  NLL, Brier score, and ECE are also lower than for \modelGru{}.

Among the post hoc logistic controls, \modelLinear{} is strongest.
\modelDistill{} leads it by .0139 BA, with a crossed reader/article interval of
$[-0.0037,0.0321]$.

\subsection{Post Hoc Ablations and Seed Sensitivity}

Table~\ref{tab:controls} averages the post hoc comparisons over the matched
random seeds.  Label smoothing does not improve BA over the GRU.  All KD
conditions improve BA at every seed, with reader-clustered intervals above
zero.

\begin{table}[t]
\caption{Post hoc matched-seed results.  Student $\Delta$ values are relative
to GRU+BCE.  Direct contrasts follow the model rows.  CIs use reader-clustered
resampling.}
\label{tab:controls}
\centering
\scriptsize
\setlength{\tabcolsep}{1.7pt}
\begin{tabular}{@{}lrrrr@{}}
\toprule
Condition & BA & $\Delta$BA & Reader 95\% CI & Seeds $>0$ \\
\midrule
GRU+BCE & .5857 & -- & -- & -- \\
$+$ label smoothing & .5852 & $-$.0005 & [$-$.0022,.0011] & 1/3 \\
$+$ gaze-only KD & .6221 & +.0364 & [.0199,.0534] & 3/3 \\
$+$ text--gaze KD & .6232 & +.0376 & [.0223,.0529] & 3/3 \\
$+$ class-shuffled KD & .6042 & +.0185 & [.0052,.0317] & 3/3 \\
\midrule
Text--gaze KD $-$ gaze-only KD & -- & +.0011 & [$-$.0030,.0052] & 1/3 \\
Text--gaze KD $-$ class-shuffled & -- & +.0190 & [.0061,.0317] & 3/3 \\
\bottomrule
\end{tabular}
\end{table}

The two teachers yield similar students under both bootstrap schemes.
Text--gaze KD exceeds class-shuffled KD at every seed, with a crossed interval
of $[.0033,.0352]$.  Class-shuffled KD remains above the GRU.

\subsection{Inference Cost}

Because distillation leaves the student unchanged, the distilled and
undistilled GRUs have identical parameter counts and inference operations.  The
serialized student state occupies 176~KiB.  Across all retained sequences,
model-only latency has a 5.08-ms median and a 16.57-ms 95th percentile.

\section{Discussion}

Teacher-derived soft targets improve the fixed gaze-only GRU without changing
its inference path.  The tested label-smoothing condition does not provide a
comparable benefit.  The gaze-summary logistic control narrows the margin, so
the evidence supports distillation for this GRU rather than superiority over
all gaze-only models.

Removing contextual text has no measurable effect on student performance.
This does not support a text-specific benefit.  Both teachers retain word-level
gaze aggregation and greater capacity than the student, so the experiment does
not separate these factors.

Within-class shuffling cuts the matched-seed gain roughly in half but does not
remove it.  Trial-level correspondence contributes to the gain.  The remainder
may reflect class-dependent soft targets or their optimization effect.

The ablations were designed after the primary result, and three seeds do not
characterize optimization variability.  The confidence intervals quantify
reader and article sampling after averaging across those seeds.  Because
reading goal is assigned between readers, a crossover study is needed to
identify within-reader effects.
The data come from one English news corpus and a laboratory eye tracker, and
all predictions use completed scanpaths.  Other domains, trackers,
self-selected goals, and partial sequences require separate evaluation.

The gaze-only student still requires calibrated coordinates and paragraph
bounds.  End-to-end device studies should include eye tracking, preprocessing,
memory, and energy, which are outside the CPU benchmark.  Any interface
adaptation should be opt-in and evaluated with users.

\section{Conclusion}

Under simultaneous reader and article holdout, teacher distillation improves a
compact gaze-only classifier at unchanged inference cost.  Within-class target
shuffling reduces but does not eliminate the gain.  Contextual text provides no
measurable improvement in student performance.

\section*{Code Availability}

{\raggedright
Code, configurations, and aggregate results are available at
\url{github.com/gohyunsu/gaze-only-reading-goal-distillation}.
\par}

\bibliographystyle{IEEEtran}
\bibliography{references}

\end{document}
